% ==============================================================
\section{Nested Structures}
\label{sec:nested-structures}
% ==============================================================

In the preceding sections, we decomposed generalized means recursively using the \emph{same} power parameter $\rho$ at every level. This ``flat'' structure arises naturally in discounted utility and expected utility when a single elasticity governs all substitution.

We now consider \emph{nested} structures: generalized means of generalized means, where different levels use different power parameters. This additional flexibility is essential for modeling situations where substitution patterns differ across groups---for instance, when goods within a category are closer substitutes than goods across categories, or when risk aversion differs from intertemporal substitution.

% --------------------------------------------------------------
\subsection{From Recursion to Nesting}
\label{subsec:recursion-to-nesting}
% --------------------------------------------------------------

Recall the recursive decomposition from \cref{sec:recursive-decomposition}. For an $n$-element mean:
\begin{equation}
\M(\bx; \bom, \rho) = \M\Big( x_1, \, \M(\bx_{2:n}; \hat{\bom}, \rho); \; \omega_1, \, \rho \Big).
\label{eq:recursion-recall}
\end{equation}
The key feature is that the \emph{same} $\rho$ appears in both the outer and inner means.

\paragraph{Breaking the Symmetry.}
Suppose instead we use a \emph{different} power parameter for the inner mean:
\begin{equation}
\M\Big( x_1, \, \M(\bx_{2:n}; \hat{\bom}, \gamma); \; \omega_1, \, \rho \Big),
\label{eq:nested-different-rho}
\end{equation}
where $\gamma \neq \rho$. This is no longer equivalent to a flat mean over all elements---it is a genuinely nested structure.

\paragraph{Two-Level Nesting.}
More generally, consider $n$ elements partitioned into groups $\mathcal{G}_1, \ldots, \mathcal{G}_K$. Define the group-level aggregates:
\begin{equation}
\bar{x}_k = \M(\bx_k; \bom_k, \gamma_k), \quad k = 1, \ldots, K,
\label{eq:group-aggregates}
\end{equation}
where $\bx_k$ contains the elements in group $k$, $\bom_k$ are the within-group weights, and $\gamma_k$ is the within-group power parameter.

The overall mean aggregates across groups:
\begin{equation}
\M^{\text{nested}} = \M(\bar{x}_1, \ldots, \bar{x}_K; \bom^{\text{outer}}, \rho),
\label{eq:overall-nested}
\end{equation}
where $\bom^{\text{outer}}$ are the across-group weights and $\rho$ is the across-group power parameter.

\begin{definition}[Two-Level Nested Mean]
\label{def:nested-mean}
A two-level nested generalized mean with groups $\mathcal{G}_1, \ldots, \mathcal{G}_K$ is:
\begin{equation}
\M^{\text{nested}}(\bx) = \M\Big( \M(\bx_1; \bom_1, \gamma_1), \ldots, \M(\bx_K; \bom_K, \gamma_K); \; \bom^{\text{outer}}, \rho \Big).
\end{equation}
The inner means use within-group parameters $\gamma_k$; the outer mean uses across-group parameter $\rho$.
\end{definition}

% --------------------------------------------------------------
\subsection{When Nesting Matters}
\label{subsec:when-nesting-matters}
% --------------------------------------------------------------

Nesting is substantive only when the power parameters differ across levels. If all $\gamma_k = \rho$, the nested structure collapses to a flat mean.

\paragraph{The Flat Case.}
Suppose $\gamma_k = \rho$ for all $k$. Then the inner means satisfy:
\begin{equation}
\M(\bx_k; \bom_k, \rho)^\rho = \sum_{i \in \mathcal{G}_k} \omega_{k,i}^{1-\rho} x_i^\rho.
\label{eq:inner-flat}
\end{equation}
Substituting into the outer mean and using the appropriate weight renormalization, we recover a single flat mean over all elements. The grouping is immaterial.

\paragraph{The Nested Case.}
When $\gamma_k \neq \rho$, the inner and outer aggregations are not interchangeable. The group structure matters: elements within a group interact through $\gamma_k$, while groups interact through $\rho$.

\begin{calloutimportant}[The Key Distinction]
\begin{itemize}[leftmargin=1.5em]
    \item \textbf{Flat structure} ($\gamma_k = \rho$): Grouping is arbitrary; any partition gives the same result.
    \item \textbf{Nested structure} ($\gamma_k \neq \rho$): Grouping is substantive; different partitions yield different values.
\end{itemize}
Nesting introduces a hierarchy where substitution within groups differs from substitution across groups.
\end{calloutimportant}

% --------------------------------------------------------------
\subsection{Example: Two Groups of Two}
\label{subsec:example-two-groups}
% --------------------------------------------------------------

Consider four elements $x_1, x_2, x_3, x_4$ partitioned into two groups: $\mathcal{G}_1 = \{1, 2\}$ and $\mathcal{G}_2 = \{3, 4\}$. Assume equal weights throughout for simplicity.

\paragraph{Flat Mean.}
With a single parameter $\rho$:
\begin{equation}
\M^{\text{flat}} = \left( \frac{1}{4} \sum_{i=1}^{4} x_i^\rho \right)^{1/\rho}.
\label{eq:flat-four}
\end{equation}

\paragraph{Nested Mean.}
With within-group parameter $\gamma$ and across-group parameter $\rho$:
\begin{align}
\bar{x}_1 &= \left( \frac{1}{2} x_1^\gamma + \frac{1}{2} x_2^\gamma \right)^{1/\gamma}, \label{eq:nested-group1} \\[4pt]
\bar{x}_2 &= \left( \frac{1}{2} x_3^\gamma + \frac{1}{2} x_4^\gamma \right)^{1/\gamma}, \label{eq:nested-group2} \\[4pt]
\M^{\text{nested}} &= \left( \frac{1}{2} \bar{x}_1^\rho + \frac{1}{2} \bar{x}_2^\rho \right)^{1/\rho}. \label{eq:nested-outer}
\end{align}

\paragraph{Comparing the Two.}
Set $x_1 = x_2 = 1$ and $x_3 = x_4 = 2$. Then:
\begin{align}
\bar{x}_1 &= 1, \quad \bar{x}_2 = 2, \\[4pt]
\M^{\text{nested}} &= \left( \frac{1}{2} \cdot 1^\rho + \frac{1}{2} \cdot 2^\rho \right)^{1/\rho} = \Mbar(1, 2; 1/2, \rho).
\end{align}
The nested mean depends only on $\rho$, not $\gamma$, because elements within each group are equal.

Now set $x_1 = 0.5$, $x_2 = 1.5$, $x_3 = 1.5$, $x_4 = 2.5$. The within-group variation activates the $\gamma$ parameter:
\begin{align}
\bar{x}_1 &= \left( \frac{1}{2} (0.5)^\gamma + \frac{1}{2} (1.5)^\gamma \right)^{1/\gamma}, \\[4pt]
\bar{x}_2 &= \left( \frac{1}{2} (1.5)^\gamma + \frac{1}{2} (2.5)^\gamma \right)^{1/\gamma}.
\end{align}
Now both $\gamma$ and $\rho$ affect the final value.

\begin{callouttip}[Exercise]
For the values $x_1 = 0.5$, $x_2 = 1.5$, $x_3 = 1.5$, $x_4 = 2.5$:
\begin{enumerate}[label=(\alph*), leftmargin=1.5em]
    \item Compute $\M^{\text{flat}}$ for $\rho = 0$ (geometric mean).
    \item Compute $\M^{\text{nested}}$ for $\gamma = 1$ (arithmetic within groups) and $\rho = 0$ (geometric across groups).
    \item Verify that the answers differ.
\end{enumerate}
\end{callouttip}

% --------------------------------------------------------------
\subsection{Economic Interpretations}
\label{subsec:economic-interpretations}
% --------------------------------------------------------------

Nested CES structures appear throughout economics. The grouping reflects economic relationships; the differing elasticities reflect differing substitution patterns.

\paragraph{Consumer Theory.}
Goods are grouped into categories: food, housing, transportation, etc. The within-category elasticity $\sigma_k$ governs substitution among goods in category $k$ (e.g., among different foods). The across-category elasticity $\sigma$ governs substitution among categories.

Typically $\sigma_k > \sigma$: it is easier to substitute chicken for beef (within food) than to substitute food for housing (across categories).

\paragraph{Production Theory.}
Inputs are grouped by type: capital, labor, energy, materials. Within labor, we might distinguish skilled and unskilled workers. The within-group elasticity governs substitution among worker types; the across-group elasticity governs substitution between labor and capital.

\paragraph{International Trade.}
Consumers choose among varieties of goods from different countries. Within a country, varieties may be close substitutes (high $\sigma_k$). Across countries, substitution may be harder (lower $\sigma$). This structure underlies the Armington model and its descendants.

\begin{calloutnote}[The Armington Aggregator]
In trade models, a common specification is:
\begin{equation}
C = \left( \sum_k \omega_k \left( \sum_{i \in k} \omega_{ki} \, c_{ki}^{(\sigma_k - 1)/\sigma_k} \right)^{\sigma_k/(\sigma_k-1) \cdot (\sigma-1)/\sigma} \right)^{\sigma/(\sigma-1)},
\end{equation}
where $k$ indexes countries and $i$ indexes varieties within country $k$. This is a nested CES with within-country elasticity $\sigma_k$ and across-country elasticity $\sigma$.
\end{calloutnote}

% --------------------------------------------------------------
\subsection{Properties of Nested Means}
\label{subsec:properties-nested}
% --------------------------------------------------------------

Nested means inherit many properties from the underlying generalized means, but some new features emerge.

\paragraph{Monotonicity and Intermediate Value.}
Nested means remain monotone in each argument and lie between the minimum and maximum elements. These properties follow from the same properties of the component means.

\paragraph{Homogeneity.}
If all component means are homogeneous of degree one (which the canonical and classical forms are), then the nested mean is also homogeneous of degree one:
\begin{equation}
\M^{\text{nested}}(\lambda \bx) = \lambda \M^{\text{nested}}(\bx) \quad \text{for all } \lambda > 0.
\label{eq:nested-homogeneity}
\end{equation}

\paragraph{Order Dependence.}
Unlike the flat case, the order of aggregation matters when $\gamma \neq \rho$. Aggregating first within groups then across groups gives a different result than aggregating in a different grouping.

\begin{proposition}[Non-Commutativity of Nesting]
\label{prop:non-commutativity}
Let $\gamma \neq \rho$. Consider two partitions of elements into groups. The nested means under the two partitions are generally different:
\begin{equation}
\M^{\text{nested}}_{\text{partition } A} \neq \M^{\text{nested}}_{\text{partition } B}.
\end{equation}
The grouping structure is a substantive modeling choice.
\end{proposition}

% --------------------------------------------------------------
\subsection{Sufficient Statistics in Nested Structures}
\label{subsec:sufficient-nested}
% --------------------------------------------------------------

The sufficient statistic property extends to nested structures, but at each level.

\paragraph{Within-Group Aggregation.}
The group mean $\bar{x}_k = \M(\bx_k; \bom_k, \gamma_k)$ is a sufficient statistic for all elements in group $k$. To compute the overall nested mean, we need only $\bar{x}_1, \ldots, \bar{x}_K$---not the individual elements.

\paragraph{Across-Group Aggregation.}
Given the group means, the outer aggregation proceeds as a standard generalized mean. The continuation value in a dynamic nested problem is itself a nested mean.

\begin{calloutnote}[Multi-Stage Budgeting]
In consumer theory, nested CES leads to multi-stage budgeting:
\begin{enumerate}[leftmargin=1.5em]
    \item The consumer allocates total expenditure across categories using the across-category price index and elasticity.
    \item Within each category, the consumer allocates the category budget across goods using the within-category price index and elasticity.
\end{enumerate}
Each stage is a standard CES problem. The category price index---a generalized mean of within-category prices---is the sufficient statistic that links the stages.
\end{calloutnote}

% --------------------------------------------------------------
\subsection{Summary}
\label{subsec:nested-summary}
% --------------------------------------------------------------

\begin{calloutimportant}[Key Results]
\begin{enumerate}[label=(\roman*), leftmargin=2em]
    \item A \textbf{nested structure} uses different power parameters at different levels of aggregation.
    
    \item When all parameters are equal ($\gamma_k = \rho$), the nested structure \textbf{collapses to a flat mean}---grouping is irrelevant.
    
    \item When parameters differ, \textbf{grouping is substantive}: different partitions yield different values.
    
    \item Nested CES is ubiquitous in economics: consumer theory (goods within categories), production (factors within types), trade (varieties within countries).
    
    \item The \textbf{sufficient statistic property} holds at each level: group means summarize within-group information for the outer aggregation.
\end{enumerate}
\end{calloutimportant}

The nested structure provides the mathematical foundation for Epstein--Zin preferences, which nest risk (aggregation over states) inside time (aggregation over periods) with different elasticities at each level.
